\subsection{Lý do chọn đề tài}

Trong bối cảnh giáo dục hiện nay, việc rèn luyện tư duy logic và khả năng giải quyết vấn đề cho học sinh tiểu học ngày càng được chú trọng. Toán tư duy không chỉ giúp các em nâng cao năng lực học thuật mà còn là nền tảng quan trọng cho các kỹ năng như phân tích, suy luận và sáng tạo. Tuy nhiên, các hình thức học truyền thống thường khô khan, thiếu tính trực quan và dễ làm giảm hứng thú của học sinh.

Bên cạnh đó, công nghệ giáo dục (EdTech) đang phát triển mạnh mẽ, mở ra cơ hội ứng dụng các phương pháp mới giúp học sinh tiếp cận kiến thức một cách sinh động và hiệu quả hơn. Việc xây dựng một trang web học toán tư duy với giao diện trực quan, hoạt động tương tác và mini-game sinh động sẽ giúp các em vừa học vừa chơi, từ đó tăng khả năng ghi nhớ và phát triển tư duy tự nhiên.

Chính vì vậy, nhóm quyết định lựa chọn đề tài “Website học toán tư duy cho học sinh tiểu học” nhằm tạo ra một công cụ học tập hiện đại, dễ sử dụng, phù hợp với lứa tuổi và có khả năng hỗ trợ việc học tập một cách hiệu quả và thú vị hơn.

\subsection{Thách thức hiện tại}

\begin{itemize}
\item \textbf{Thiếu môi trường học tập trực quan và tương tác:} Phần lớn tài liệu toán tư duy hiện nay ở dạng sách in hoặc bài tập truyền thống, khiến học sinh khó hình dung, đặc biệt với các bài đòi hỏi suy luận hoặc thao tác trực tiếp.
\item \textbf{Khó duy trì hứng thú học tập:} Nội dung học không được game hóa hoặc thiết kế hấp dẫn khiến học sinh dễ chán nản, không duy trì được sự tập trung trong thời gian dài.

\item \textbf{Thiếu hệ thống theo dõi tiến trình học tập:} Giáo viên và phụ huynh khó nắm bắt được mức độ hiểu bài của học sinh, tiến trình học hoặc điểm mạnh – điểm yếu để có hướng hỗ trợ kịp thời.

\item \textbf{Chưa có công cụ học toán tư duy phù hợp độ tuổi:} Nhiều nền tảng hiện nay thiết kế cho lứa tuổi lớn hơn, sử dụng thuật ngữ phức tạp hoặc giao diện khó sử dụng đối với học sinh tiểu học.

\item \textbf{Hạn chế trong việc cá nhân hóa trải nghiệm học:} Mỗi học sinh có khả năng tiếp thu khác nhau, nhưng các tài liệu hiện tại chưa đáp ứng tốt việc tùy chỉnh độ khó hoặc cung cấp gợi ý phù hợp năng lực từng em.
\end{itemize}

